% from ega3-6-fr.tex
\oldpage[III]{171}
On a alors des isomorphismes canoniques fonctoriels
\[
  \mathfrak{Tor}_n^{S'}
  (f',1_{S'};\mathscr P'_\bullet,\mathscr F')
  \xrightarrow{\sim}
  \mathscr H^{-n}(f,\mathscr P_\bullet)
  \otimes_{\mathscr O_S}\mathscr F'
  \tag{6.9.9.1}\label{III.6.9.9.1.fr}
\]
où
$\mathscr P'_\bullet=\mathscr P_\bullet\otimes_{\mathscr O_S}
\mathscr O_{S'}$; en particulier, pour $\mathscr F'=\mathscr O_{S'}$, on a
des isomorphismes canoniques fonctoriels
\[
  \mathscr H^n(f',\mathscr P'_\bullet)
  \xrightarrow{\sim}
  \mathscr H^n(f,\mathscr P_\bullet)
  \otimes_{\mathscr O_S}\mathscr O_{S'}.
  \tag{6.9.9.2}\label{III.6.9.9.2.fr}
\]
C'est le cas particulier $m=1$ de (6.9.6). Plus particulièrement :
\end{corollary}

\begin{corollary}[6.9.10]
Soient $S$ un préschéma, $f:X\to Y$ un $S$-morphisme séparé et
quasi-compact de $S$-préschémas, $\mathscr P_\bullet$ un complexe limité
inférieurement, formé de $\mathscr O_X$-Modules quasi-cohérents, plats sur
$S$. On suppose en outre que les $\mathscr O_Y$-Modules
$\mathscr H^n(f,\mathscr P_\bullet)$ soient plats sur $S$. Pour tout $s\in S$,
notons $X_s$ et $Y_s$ les fibres
$X\otimes_S k(s)$, $Y\otimes_S k(s)$,
$f_s:X_s\to Y_s$ le morphisme $f\times_S1$,
$\mathscr P_\bullet^s$ le complexe
$\mathscr P_\bullet\otimes_{\mathscr O_S}k(s)$ de
$\mathscr O_{X_s}$-Modules. On a alors des isomorphismes canoniques
fonctoriels
\[
  \mathscr H^n(f_s,\mathscr P_\bullet^s)
  \xrightarrow{\sim}
  \mathscr H^n(f,\mathscr P_\bullet)\otimes_{\mathscr O_S}k(s).
  \tag{6.9.10.1}\label{III.6.9.10.1.fr}
\]
\end{corollary}

On a donc, moyennant des hypothèses de platitude convenables, un cas où la
formation des foncteurs dérivés $R^nf_*(\mathscr F)$ « commute au passage aux
fibres », que nous retrouverons par une autre méthode au §~7.

\subsection{Structure locale de certains foncteurs cohomologiques.}

\begin{proposition}[6.10.1]
Soient $S=\operatorname{Spec}(A)$ un schéma affine,
$Y^{(i)}$ ($1\leq i\leq n$) une famille finie de $S$-schémas affines, plats
sur $S$; pour chaque $i$, soit $f_i:X^{(i)}\to Y^{(i)}$ un $S$-morphisme
séparé et quasi-compact, et soit $\mathscr P_\bullet^{(i)}$ un complexe de
$\mathscr O_{X^{(i)}}$-Modules quasi-cohérents limité inférieurement. Soit
$Y$ le produit des $S$-schémas $Y^{(i)}$. Il existe un complexe
$\mathscr R_\bullet$ de $\mathscr O_Y$-Modules quasi-cohérents et plats sur
$S$, ayant la propriété suivante : pour tout $S$-schéma affine $S'$ et tout
complexe de $\mathscr O_{S'}$-Modules quasi-cohérents
$\mathscr Q'_\bullet$ limité inférieurement, il y a un isomorphisme
\[
  \mathbf{Tor}_\bullet^S
  (f_1,\ldots,f_n,1_{S'};
   \mathscr P_\bullet^{(1)},\ldots,\mathscr P_\bullet^{(n)},
   \mathscr Q'_\bullet)
  \xrightarrow{\sim}
  \mathscr H_\bullet
  (\mathscr R_\bullet\otimes_{\mathscr O_S}\mathscr Q'_\bullet)
  \tag{6.10.1.1}\label{III.6.10.1.1.fr}
\]
qui est un isomorphisme de $\partial$-foncteurs en $\mathscr Q'_\bullet$. En
outre, pour tout $S$-morphisme $u:S''\to S'$ de $S$-schémas affines, le
diagramme
\[
\xymatrix@C=12pt@R=30pt{
  {\begin{gathered}
    \mathbf{Tor}_\bullet^S
    (f_1,\ldots,f_n,1_{S'};\\
    \mathscr P_\bullet^{(1)},\ldots,\mathscr P_\bullet^{(n)},
    \mathscr Q'_\bullet)
  \end{gathered}}
  \ar[r]^-{\sim}\ar[d]
  &
  {\mathscr H_\bullet
    (\mathscr R_\bullet\otimes_{\mathscr O_S}\mathscr Q'_\bullet)}
  \ar[d]
  \\
  {\begin{gathered}
    \mathbf{Tor}_\bullet^S
    (f_1,\ldots,f_n,1_{S''};\\
    \mathscr P_\bullet^{(1)},\ldots,\mathscr P_\bullet^{(n)},
    u^*(\mathscr Q'_\bullet))
  \end{gathered}}
  \ar[r]^-{\sim}
  &
  {\mathscr H_\bullet
    (\mathscr R_\bullet\otimes_{\mathscr O_S}u^*(\mathscr Q'_\bullet))}
}
\tag{6.10.1.2}\label{III.6.10.1.2.fr}
\]
(où les flèches verticales sont les $(1_Y\times_Su)$-morphismes canoniques
définis dans (6.7.10)) est commutatif.
\end{proposition}

\begin{proof}
